% Capstone 5: Lenstra's ECM. Port of «Мини-курсовая - алгоритм факторизации Ленстры».
% The performance numbers come from the generated requirements.tex (make bench).
\documentclass{ca-capstone}
\capstonerefs{sagemath,crandall-pomerance,hps}
\begin{document}
% Generated by `make bench CAPSTONE=lenstra-ecm` (taskgen.capstone report) from
% bench/results.json: reference time x slack on the reference machine.
% Do not edit: change bench/config.yaml and re-run the benchmark.
\newcommand{\ReqSlack}{10}
\newcommand{\ReqBitsMin}{50}
\newcommand{\ReqBitsMax}{100}
\newcommand{\ReqSeconds}{50}

\capstonetitle{Факторизация алгоритмом Ленстры}

\goal
Требуется реализовать алгоритм Ленстры разложения целого числа на множители.

\functionality
\begin{enumerate}
  \item Порождение кривой $(\mathrm{mod}\ N)$ и точки на ней\footnote{как это сделать?}
  \item Нахождение нетривиального фактора входного целого числа~$N$
  \item Разложение входного целого числа~$N$ на простые множители
\end{enumerate}

\requirements
Требования к реализации:
\begin{enumerate}
  \item корректность решения задачи факторизации
  \item приемлемая скорость работы (см. комментарий ниже)
\end{enumerate}

\textbf{Корректность работы} можно проверять, самостоятельно генерируя задачи.

В данной работе важна не скорость, а корректность реализации и соответствие требованиям.
Однако тестовые примеры с числами размеров примерно \textbf{\ReqBitsMin--\ReqBitsMax~бит}
должны выполняться за приемлемое время "--- несколько десятков секунд (не более
\ReqSeconds~с).

\details
Так как цель работы "--- решение задачи разложения на множители, в своей реализации
\textbf{можно пользоваться библиотеками} для работы с эллиптическими кривыми, а также пакетами
линейной алгебры.

Рекомендуется использовать язык \emph{Python} и систему компьютерной алгебры
\emph{SageMath}~\cite{sagemath}.

Желательно сохранять промежуточные результаты алгоритма Ленстры, такие как уравнение
сгенерированной кривой и координаты точки~$P$, координаты вычисляемых точек~$n!\,P$ и проч.

Платформо-независимый, \textbf{переносимый код} писать \textbf{не требуется}, но во время сдачи
должна быть возможность проверить работу.

\schedule
Работа выполняется \textbf{индивидуально}. Срок выполнения "--- \textbf{месяц}.

Работа должна быть ваша: реализация должна быть \textbf{разработана и написана вами}, а не
собрана из кусков чужих работ.

Желательно хранить код в \emph{git}-репозитории на \emph{github}, \emph{bitbucket} или подобных
платформах. \textbf{Документация не требуется}.

\remarks
\begin{enumerate}
  \item Существенным плюсом работы будет реализация следующих оптимизаций базового алгоритма:
    \begin{itemize}
      \item специальная параметризация для порождения кривых
      \item выбор кривых с порядком, делящимся на 12 или 16
    \end{itemize}
    см. п.~7.4.2, теорему~7.4.3 и алгоритм~7.4.4 в~\cite{crandall-pomerance}
\end{enumerate}

\references
\end{document}
